% !TeX root = ../papers/nc-p3.tex
\section{有限辺の平方根の接空間と階数十のインスタントン}
\label{sec:finite-square-tangent-instantons}
四辺有限な標準有限Frobenius窓において、一つの辺の文字空間を $U,V\subset D$、関係を $R\subset U\otimes V$、
$K=R^\perp=(UV)^*$ と書く。ここで $D$ は長さ20の点グラフの関数環である。
定理\ref{thm:finite-square-chart}の平方根の閉スキーム
$\mathcal F_Q=\{L\in\Gr(4,D):L^2\subset Q\}$、$Q=U^2$ の既知の根 $U$ に注目する。
以下は辺が非被約な場合も含む。
\begin{theorem}\label{thm:finite-square-tangent-instantons}
辺の完全跡のスカラーを選べば、$\mathcal F_Q$ の $U$ での接空間は
\[
 \bigwedge\nolimits^2 K\ \cap  \left(\bigwedge\nolimits^2 U^*\otimes\operatorname{Sym}^2 V^*\right)
 \quad\subset\quad \bigwedge\nolimits^2(U\otimes V)^*
\]
に同定される。交代形式の非零な階数は $6,8,10$ のいずれかである。
階数10の方向は $\mathbb P(U)$ 上の階数2、電荷4の数学的インスタントン束を与え、
そのモナドの線型写像からもとの辺 $R$ を復元できる。
階数6,8の方向では、その根基は必ず可分解テンソルを含む。
したがって二つの交代化行列の線型代数だけで、各向きで正規化した根が被約孤立点かどうかを判定できる。
\end{theorem}
\begin{proof}[接空間の同定]
$U$ から $D/U$ への無限小移動 $f$ の条件は $u f(u')+u'f(u)=0\bmod Q$。
辺の跡で $Q^\perp=V^2$、$U^\perp=UV^2$ である。後者では包含は跡の $(2,2)$ 消滅による。
$(1,2)$ の関係は $R\otimes V$ の24本で全てかつ独立である。六本の完全交叉のイデアルの
Koszul分解を $(1,2)$ でひねると、定義式からの6組の一次式以外に全次数0以下の寄与はない
（両因子とも最高次コホモロジーを持つ項は外次数6、全次数はイデアルの分解で1）。
従って $\dim UV^2=40-24=16$ で等号となる。
接写像の双対の余核は
\[
 (U\otimes UV^2)/\langle u\otimes u'q+u'\otimes uq:\ q\in V^2\rangle
   = (\bigwedge^2 U\otimes\operatorname{Sym}^2 V)/
      \langle(u\wedge r_U)\otimes (v r_V):u\in U,v\in V,r\in R\rangle .
\]
最後は $r=\sum r_U\otimes r_V$ に総和を取る記法である。
この空間の双対は、$U$ 側で交代、$V$ 側で対称かつ $R$ と任意の $u\otimes v$ の組で
零となる形式である。これが表示した交叉を与える。接空間が零であることと
有限型スキームのその局所環が剰余体であることは中山の補題で同値である。
\end{proof}

\begin{lemma}\label{lem:finite-four-forms-evaluate-three}
有限辺の $K$ 内の任意の四次元空間 $W$ は一般の $u\in U$ で $V^*$ に三次元以上の像を持つ。
$U,V$ を逆にしても同じである。
\end{lemma}
\begin{proof}
$\mathbb P(K)$ 内の階数一の点は有限である。さらに $K$ は、片側の支持が二次元の
四次元行列空間を含まない。そのような空間の射影三平面は
$\mathbb P^1\times\mathbb P^3\subset\mathbb P^7$ に固有交叉し、その交叉長は次数4である。
これは点グラフの片側の直線切断にスキームとして入り、系\ref{cor:finite-no-four-on-line}に反する。

全 $u$ で評価階数が2以下と仮定する。$Mu=0$ となる
$([M],[u])\in\mathbb P(W)\times\mathbb P(U)$ の接続は $u$ 側の各繊維が
次元1以上で、全次元は4以上である。$M$ 側からは階数3の繊維は一点、階数2では直線、
有限個の階数1では平面である。従って $W$ の一般元が階数2でなければならず、
階数1の有限集合を除く開集合で定階数2である。

この開集合で $U\otimes\mathcal O\twoheadrightarrow E\hookrightarrow V^*\otimes\mathcal O(1)$
を普遍評価の像で分解する。$E$ と $E^*(1)$ は階数2で大域生成される。
有限集合を除く $\mathbb P^3$ のPicard群は $\mathbb Z$ なので $\det E=\mathcal O(t)$、
$0\le t\le2$ となる。$t=0,2$ のいずれかなら、行または列二平面のPl\"ucker写像が定数となり、
排除済みの片側二次元支持になる。$t=1$ なら二次元核のGrassmann写像は一次式で与えられる。
その式は全 $\mathbb P(W)$ に延長し、共通零点は有限である。
その像の閉包は $\Gr(2,U)$ に含まれる射影線型空間で、次元は2以上。
Klein二次曲面内の極大線型空間は平面であり、全ての直線が一点を通る型か、
一つの平面内の直線全体の型に限る。これはPl\"uckerの偏極から任意の二直線が交わることでも分かる。
核の合併は $\mathbb P(U)$ を覆う（有限個の階数1の核を加えても稠密性は変わらない）ので
後者は不可能である。よって全 $M$ に共通の核ベクトル $u_0$ がある。
$\bar U=U/ku_0$ とすると、動く核直線は全射な線型写像 $L:W\to\bar U$ による。
$M L(M)=0$ は全 $M$ での二次恒等式である。$0\ne M_0\in\ker L$ を取り偏極すれば
$M_0 L(M)=0$ が全 $M$ で成立する。$L$ が全射で $M_0u_0=0$ なので $M_0=0$、矛盾である。
\end{proof}

\begin{proof}[交代階数とモナド]
非零な接方向の根基を割ったシンプレクティック空間を $H$ とし、$H^*\subset K$ を考える。
各 $u$ で $uV$ の $H$ における像は等方的で、次元は $\dim H/2$ 以下である。
階数4なら $H^*$ の四本の形式の $u$ での評価が全て二次元以下になり、補題に反する。
階数2なら $H^*$ の鉛筆の形式が各 $u$ で一次元以下に評価される。
従って鉛筆の行列の核が $\mathbb P(U)$ を覆うが、一般の行列の階数が2以上なら
これらの核の合併は高々二次元であり、有限個の階数1を加えても覆えない。
鉛筆全体が階数1となるのも有限辺に反する。

$\mathbb P(U)$ 上に
\[
  V\otimes\mathcal O(-1)\xrightarrow{A}H\otimes\mathcal O
       \xrightarrow{A^*\omega} V^*\otimes\mathcal O(1)
\]
を作る。合成は等方性で零である。根基が純テンソルを避ければ、$A$ は全点で階数4となる。
階数10の場合は根基がちょうど $R$ であり、有限辺には純関係がないのでこの場合に入る。
中間コホモロジー束 $F$ は階数2、$c(F)=(1-h^2)^{-4}$、
従って $c_1=0,c_2=4$ である。$H^0(F)=0$ は $H^*\hookrightarrow (U\otimes V)^*$ と
$H^1(\mathcal O(-1))=0$ から、$H^1(F(-2))=0$ はモナドの捻りから従う。
階数2かつ $c_1=0$ の正規化された束では $H^0(F)=0$ は傾き安定性を与え、インスタントンの条件を満たす。
線型作用 $U\otimes V\to H$ の核で辺を復元する。
階数6なら四次元等方部分空間はなく、階数8で全点階数4なら中間コホモロジーが零なのに
全Chern類が $1+4h^2$ となる。従っていずれも根基に純テンソルがある。
\end{proof}

\begin{corollary}\label{cor:opposite-square-root-collision}
離散付値環上で四辺有限の標準窓が動き、一般繊維では反対二位置が単元倍で一致せず、
特殊繊維では一致するとする。特殊繊維のその向きの辺には上記の階数6以上の交代形式が存在する。
\end{corollary}
\begin{proof}
完備化と有限拡大後に二位置の二乗を揃える。実際有限平坦な相対点環の各文字直線束は
半局所環上で自明化でき、二つの辺の完全跡の比は単元でその平方根を取れる。
特殊繊維で根そのものも一致するよう符号を選べる。二乗空間を保つ単元はスカラーだけである。
実際単元 $s\in U$ に対し $g U^2\subset U^2$ なら、$g s^2$ と $g s u$ を二次式と読み
$(g s^2)u=(g s u)s$ を三次式の単射性で持ち上げると $g s^2=s v$ となる。
さらに $g u^2\in U^2$ より $s\mid v u^2$ なので $v$ は $s$ の定数倍である。
平方1の単元は局所因子ごとの符号で、冪等元とともに完備環へ持ち上がる。
二つのGrassmann切断の差が初めて現れる付値を取れば、平方根の方程式の差の先頭項は
特殊繊維の相対接方程式を満たす。差は非零であり定理を適用できる。
\end{proof}

\paragraph{有限辺における階数十の厳密例。}
\path{computations/verify_finite_square_tangent_instantons.py} は有理数上で
$10\times4$ の線型行列 $A(x,y,z,w)$ を構成する。上四行は $x+a_i(z,w)$ の対角、
次の四行は対称行列 $yI+B(z,w)$、下二行は一次式 $f_i,g_i$ である。
$(a_i-a_j)B_{ij}=f_jg_i-f_i g_j$ により $A^*\omega A=0$ を満たす。
係数写像の階数は10であり、その核の六関係の両射影で一般平面上の十五個の最大小行列式が
全四次式を張る（コードで一つの平面を各々検査）。故に両射影は有限で、六双一次式は
長さ20の完全交叉を定める。各点でどちらの評価も階数3であり両射影は閉埋込みである。
この例に接空間の同定だけを適用するのに四窓への延長は不要である。
二つの向きの接次元は $1,0$、反転三次交叉の余階数は $4,0$ である。
従って有限な第一辺での接方向は実在し、反対側の反転三次交叉が零のため両位置を反転する標準窓へは延びない。
